\begin{figure}[tb]
\centering
\begin{tikzpicture}[>=Latex,font=\small,scale=1.05]
\draw[Purple,very thick] plot[smooth] coordinates {(-3,-1.1)(-1.7,-.2)(0,.1)(1.7,-.15)(3,-1.0)};
\node[Purple] at (2.7,-1.35){\(M=M_0\)};
\coordinate (P) at (0,.1);
\fill[BrickRed] (P) circle (2pt);
\draw[->,MidnightBlue,very thick] (P)--(0,2.0) node[above]{\(\nabla M\)};
\draw[->,ForestGreen!70!black,very thick] (P)--(1.1,1.45) node[right]{\(\matr G^{-1}\nabla M\)};
\draw[->,Orange,very thick] (P)--(2.0,-.18) node[right]{projected voltage direction};
\draw[dashed,gray] (-2.4,.6)--(2.5,.6);
\node[draw,rounded corners,fill=Orange!8,align=left] at (6.2,.55)
 {normal: change torque\\tangent 1: manage voltage\\tangent 2: remaining EESM allocation};
\end{tikzpicture}
\caption{A local task frame built from gradients. A physical metric changes
which direction is cheapest, so Euclidean orthogonality is not automatic.}
\label{fig:local-gradient-frame}
\end{figure}
